% Nouveau EGA V (cote 25): archive cote 25, site page 11 (PDF page 12).
% typescript. Draft compared: none.
% AI-transcribed by Claude Opus 5.5 (claude-opus-5-5, adaptive thinking, effort high), Message Batches API, batch ega25b (msgbatch_01A7n5YuXubUPNK1zsepi8RG), 2026-11-13; not yet checked by eye.
\scanpage{25/11}
\mspage{8}

\uline{Proposition }2.8.  Soit F cohérent sur X, et \unsure{soit} $\over{x}{y}\in Y$ , x son image
dans X. \xout{\unsure{Pxxxxquexx\&xxxxxxx}}  Soit P\inserted{(M)} une des propriétés suivantes pour
un module de type fini\inserted{M} sur un anneau local noethérien A:

(i)   coprof $M\leqslant n$   (réf).

(ii)  M satisfait$(S_k)$  (réf).

(iii)  M est \unsure{Cohen-Macaulay|CMahn-Macaulay}.

(iv)  M est réduit(\hand{\unsure{réf|vif}}).

(v)   M est intègre(\hand{\unsure{réf|vif}}).

Alors, pour que $\over{F}{G}_{\hand{\nu},y}$ satisfasse la propriété P, il f et s que $F_x$ la
satisfasse.

Cela résulte aussitôt des résultats du par.6 , compte tenu que
\xout{X} $Y_{\hand{\nu}}\longrightarrow X$  est un morphisme régulier, \unsure{donc}  $\underline{O}_{X,\unsure{x}}\longrightarrow \underline{O}_{\unsure{Y}_{\hand{\nu}},y}$ est
régulier.  Compte tenu de 2.3. on obtient donc:

\uline{Corollaire} 2.\unsure{9}.  Avec les notations de 2.8. , soit \deleted{$\unsure{E_P(F)}$}\inserted{Z} l'ensemble
des $x\in X$ tels que $P(F_x)$ ne soit pas satisfaite. \unsure{Alors} pour que
$G_{\hand{\nu}}$ satisfasse à P en tous les points, il f et s que $f(\unsure{Z})$ soit une
partie finie de P, ou encore que \xout{f(Z)} dim $\overline{f(Z)}$ = 0.

En effet, 2.8. nous dit que $\unsure{h}^{-1}(Z)$ est l'ensemble P-singulier
de $G_{\hand{\nu}}$ , et il est vide sss f(Z) est fini en vertu de 2.3. (NB h désig
\margin{|} ne le morphisme $Y_{\hand{\nu}}\longrightarrow X$; je m'aperçois à l'instant que la lettre P
\margin{|} dans 2.8. fait double emploi ).

\margin{|} \uline{Corollaire} 2.\unsure{10}.  Condition pour que $Y_{\hand{\nu}}$ soit réduit resp. loc intègre.

\deleted{On fera attention que 2.10. \inserted[typed]{et 2.11. ci-dessous} ne dit rien sur les propriétés
``géométriquement réduit'' et ``géométriquement intègre'', les propriétés
envisagées dans 2.8. étant en effet ``absolues'' . Nous \unsure{allons}}

\margin{11}
